
% Cheyne Homberger, Intro to LaTeX 2012  -  cheyne.homberger@gmail.com
\documentclass[11pt]{article}

% fonts 
\usepackage{palatino, mathpazo}     % I like this font family better than the standard


% for hyperlinks: use command \url{...}
\usepackage{hyperref}


% page setup 1 inch automatically added to margins. \oddsidemargin is the right
% margin on odd pages (even pages default to same value). The left margins are
% given by rightmargin = 8.5 - textwidth - leftmargin.
% \setlength{\oddsidemargin}{.25in} 
% \setlength{\topmargin}{-.5in}
% \setlength{\textwidth}{6in} 
% \setlength{\textheight}{8in}

% OR just use the geometry package to set all margins
\usepackage[tmargin=1in,bmargin=2in,lmargin=1in,rmargin=1in]{geometry}
\setlength{\headheight}{1in}

% \usepackage{pgfpages}              % uncomment these and compile
% \pgfpagesuselayout{2 on 1}         % to make a handout


% for doublespace, use linespread 1.6. for half space, 1.3  
\linespread{1}


% headers and footers
\usepackage{fancyhdr}  % gives greater control over headers/footers
\usepackage{graphicx}  % allows insertion of graphics
% \lhead{\includegraphics[height=1.3cm]{icubedlogo}} % old headers...
% \rhead{\includegraphics[height=1.2cm]{gmalogo}} 
\lhead{Intro to \LaTeX{}}
\rhead{{University of Florida}}
\cfoot{}            % removes default page number from center footer
\pagestyle{fancy}   % tells all pages to use the fancy style


% theorems and definitions
\usepackage{amsmath,amsthm,amssymb} 
\newtheorem{thm}{Theorem}
\newtheorem{prop}{Proposition} 
\newtheorem{defn}{Definition}   


% extra commands to save on typing...
\newcommand{\Lt}{\LaTeX{} } \newcommand{\bs}{\textbackslash}

% to create to commands which accept arguments, put the number of arguments in
% square brackets. This command centers text with \ct{this text will be
% centered}. 
\newcommand{\ct}[1]{\begin{center} {#1} \end{center}}



% ================================================================= %
\begin{document}

\title{A Very Short Introduction to \Lt} 
\date{}       % puts current date by default
\author{}     % 
\maketitle
\thispagestyle{fancy} % by default, TeX removes headers from first page

\section*{Getting Started}  % the star tells latex not to number the section
  \LaTeX{}, at its core, is a programming language used to specify the layout of
  a page. Once the code is written, it can then be compiled to produce a pdf,
  dvi, or ps document. There are a number of applications which exist
  specifically to help write and compile \Lt documents, but code can be written
  on any text editor and compiled directly. 
  
  The first step to compiling your own code is installing a \Lt distribution,
  which includes the actual programs used to compile your code into a .pdf
  document (or a .dvi, or a .ps document).  With this is installed, you can
  write in any text editor (such as notepad, vim, emacs, gedit, etc.), and build
  it into a pdf with the command \verb+pdflatex+. Even though all you need is
  the distribution and a basic editor, it is typically easier to learn on an
  editor specifically suited to \Lt.
  
  The installation process can vary depending on your operating system, but
  some of the more popular options are described below:

  \subsection*{Windows} 
    Miktex is an excellent \Lt distrubution for the
    Microsoft Windows operating systems, which is easily installed and manages
    add-on packages for you as needed. It can be downloaded from here:
    \url{http://miktex.org/}.

    For the editor, TeXnic Center (\url{http://www.texniccenter.org/}) is a \Lt
    development suite which includes references, syntax highlighing, and
    preview options.  When starting up TeXnic Center for the first time, it may
    ask you to tell it where your Miktex files are. This will typically be
    \verb+C:\Program files\MiKTeX\miktex\bin +.

  \subsection*{Mac OS X} 
    MacTeX (\url{http://www.tug.org/mactex/2011/}) is the
    most popular \Lt distribution on OS X, and it includes an excellent editor
    (TeXshop). MacTeX is a large program which is easy to set up and full of
    useful features.

  \subsection*{Linux} 
    The most popular distribution for Linux is TeXLive, which
    can be installed in a variety of ways depending on your specific Linux
    distribution (in Ubuntu: \verb+sudo apt-get install texlive+). 
    Kile is a very popular and easy to set up \Lt editor which includes many 
    useful extra features. There are also a number of plugins for editors such
    as vim or emacs. % vim, of course, being ideal

  Note that there are numerous online resources for learning \Lt (one of the
  most comprehensive and well written can be found here:
  \url{http://tobi.oetiker.ch/lshort/lshort.pdf}). If you run into trouble
  while installing, visit your particular distribution's website or try
  googling you're problem. 

\section*{Compiling Your First Document}
  Once you have a \Lt distribution installed, you're ready to compile your
  first document. Fire up any text editor, and write the following: 
  % the verbatim environment prints text exactly as written 
  \begin{verbatim}
  \documentclass{article}   

  \begin{document}

  hello world! $e^{i \pi} + 1 = 0$

  \end{document} \end{verbatim} 

  Then save this file as \verb+hello.tex+. If you're
  using an editor built for \Lt, there should be a compile button which will
  automatically turn your code into a pdf. To compile manually, pull up a
  command line, navigate to your file, and run the command 
  \verb+pdflatex hello.tex+.  Either way, you should end up with a pdf
  containing only the line:
  
  \vspace{.5cm}

  hello world! $e^{i \pi} + 1 = 0$

  \vspace{.5cm}


\section*{\Lt Syntax}
  Looking at the above example, a few details stand out. The source code is a
  mix of text and formatting commands, and the math is written in between \$
  signs. A command in \Lt typically is typically begun with a backslash, and
  arguments are passed in brackets (with options in square brackets) like so:
  \verb+\command[option]{argument1}{argument2}+.  For example, 
  \verb+\frac{a}{b}+ produces $\frac{a}{b}$. 

  The commands before the \verb+\begin{document}+ statement are known as the
  preamble and is where you place formatting options to change the appearance
  of your document. For example, the \verb+\documentclass{article}+ command
  tells \Lt that we want to use the built in class `article'.  A class sets
  basic formatting options, such as fonts, margins, formatting for section
  headings, additional packages, spacing, etc.  The `article' class is very
  flexible and easily extensible. Other options include `book', `report', and
  `beamer' (for presentations).

  The part of the document between the \verb+\begin{document}+ and
  \verb+\end{document}+ lines is the body. This is where the actual content
  goes which will make up the document. Text is divided up into regular mode
  and math mode, and math can be either inline or in display mode. By default,
  text is assumed to be non-math. To format a mathematical expression, surround
  it by dollar signs or by \verb+\( ... \)+  to fit it in the current line, or
  surround it by double dollar signs or \verb+\[ ... \]+ to display it in its
  own line. For example, 

  \verb+$e^x = \sum_{k \geq 0} \frac{x^k}{k!}$+ will compile to $e^x = \sum_{k
  \geq 0} \frac{x^k}{k!}$, while
  \\ \verb+\[ e^x = \sum_{k \geq 0} \frac{x^k}{k!} \]+ 
  will come through as \[ e^x = \sum_{k \geq 0}
  \frac{x^k}{k!}.\]

  Another important aspect of \Lt is its handling of whitespace. The code:
  \begin{verbatim}\LaTeX         treats multiple      spaces as   
  just   a single  space,  and multiple 
             blank        lines as    a   
  single     one .\end{verbatim} 
  is typeset as \ct{\Lt treats multiple spaces as just a single
  space, and multiple blank lines as a single one.} A line break is indicated
  by a pair of backslashes, and a paragraph break by at least one blank line
  between blocks of text.


\section*{Adding packages} 
  While the 'article' class is useful for most situations, and is easily
  extended to fit many situations, there are other classes specially suited for
  different tasks (notably: beamer, for presentations). Changing the formatting
  of an entire document (for example, for submission to a journal) is often as
  simple as changing the class.

  Where the class defines some of the global specifications for the document,
  packages generally add functionality or environments to a document. For
  example, this document (whose source code can be found at
  \url{http://www.math.ufl.edu/~cheyne42/texintro}), uses the package `fancyhdr' to
  make nice and customizable headers and footers, and `hyperref' to embed
  hyperlinks.

  For another example, the packages `amsmath', `amsthm', `amssymb' provide
  additional mathematical symbols and theorem-like environments, which can be
  specified in the preamble.  To see how this works, try compiling the
  following code.  Anything written after a $\%$ sign is a comment and won't
  affect the code, which can be useful for explaining what your code does (to
  youself or someone else).  \pagebreak

  \begin{verbatim} 
  \documentclass{article}

  \usepackage{amsmath,amssymb,amsthm}

  \newtheorem{defn}{Definition}  % I called them defn and thm
  \newtheorem{thm}{Theorem}      % for shorthand

  \begin{document}     
  
  \begin{defn} 
  Let $n,k \in \mathbb{Z}^+$ with $n \geq k$. Then define 
  $$ \binom{n}{k} = \frac{n!}{k!(n-k)!}.$$ 
  \end{defn}

  \begin{thm}[Binomial Theorem] 
  For all $x \in \mathbb{R}$ and $n \in \mathbb{Z}^+$, 
  $$ (x+1)^n = \sum_{k = 0}^n \binom{n}{k} x^k.$$   
  \end{thm}

  \begin{proof} Left as exercise.  \end{proof}

  \end{document} \end{verbatim}

  You should end up with something like:

  \begin{defn} 
  Let $n,k \in \mathbb{Z}^+$ with $n \geq k$. Then define
  $$\binom{n}{k} = \frac{n!}{k!(n-k)!}.$$ 
  \end{defn}

  \begin{thm}[Binomial Theorem] 
  For all $x \in \mathbb{R}$ and $n \in \mathbb{Z}^+$, 
  $$ (x+1)^n = \sum_{k = 0}^n \binom{n}{k} x^k.$$ 
  \end{thm}

  \begin{proof} Left as exercise.  \end{proof}

  More information about various packages and useful tips can be found at \\
  \url{http://en.wikibooks.org/wiki/LaTeX/}, and a more complete (and better
  written) introduction can be found at 
  \url{http://tobi.oetiker.ch/lshort/lshort.pdf}.

\end{document}
